\begin{explanationbox}

	\textbf{\textcolor{CExplanation}{一句话直觉}}

	二倍角公式揭示了角度加倍时三角函数值之间的乘积关系，是三角函数恒等变换的核心工具。

	\medskip
	\textbf{\textcolor{CExplanation}{核心拆解}}
	\begin{description}[style=nextline, leftmargin=2.8em, labelsep=0.5em, itemsep=0.45em]
		\item[\textbf{要点一（正弦余弦乘积）：}] $\sin 2x = 2\sin x \cos x$，体现正弦与余弦的对称组合。
		\item[\textbf{要点二（余弦等价形式）：}] $\cos 2x$ 有三种等写形式：$\cos^2 x - \sin^2 x$、$1-2\sin^2 x$、$2\cos^2 x-1$，不同场合选择使用。
		\item[\textbf{要点三（正切分式结构）：}] $\tan 2x = \frac{2\tan x}{1-\tan^2 x}$，使用时需保证 $\cos x\ne0$ 且 $1-\tan^2x\ne0$。
	\end{description}

	\medskip
	\textbf{\textcolor{CExplanation}{几何本质（单位圆旋转）}}

	从单位圆看，角度 $2x$ 的坐标可由角度 $x$ 的旋转复合得到，也可从和角公式直接推出二倍角公式。

	\medskip
	\textbf{\textcolor{CExplanation}{代数意义（和角特例）}}

	令和角公式 $\sin(\alpha+\beta)$ 与 $\cos(\alpha+\beta)$ 中 $\alpha=\beta=x$ 即得二倍角公式；再结合商数关系，可推出正切二倍角公式。

	\medskip
	\textbf{\textcolor{CExplanation}{考点价值（恒等变形核心）}}
	\begin{itemize}[leftmargin=*, itemsep=0.5em, topsep=0.4em]
		\item \textbf{考法一（化简求值）：} 已知 $\sin x,\cos x$，或已知其中一个三角函数值并给出象限、符号等附加信息时，求二倍角函数值。
		\item \textbf{考法二（降幂升幂）：} 利用余弦二倍角变形实现幂次的升降，为化简、证明和求值提供便利。
	\end{itemize}

	\medskip
	\textbf{\textcolor{CExplanation}{顿悟点}}

	\textbf{二倍角公式不是新知识，而是和角公式的推论，记牢和角公式可随时导出。}

	\medskip
	\textbf{\textcolor{CExplanation}{使用场景}}
	\begin{itemize}[leftmargin=*, itemsep=0.5em, topsep=0.4em]
		\item \textbf{场景一（直接求值）：} 已知 $\sin x,\cos x$，或已知其中一个三角函数值并给出象限信息时，求 $\sin 2x$、$\cos 2x$、$\tan 2x$。
		\item \textbf{场景二（恒等证明）：} 化简三角恒等式时，将二倍角展开为单角形式。
		\item \textbf{场景三（方程求解）：} 解含 $\sin 2x$、$\cos 2x$ 的三角方程，利用公式化为单角方程。
	\end{itemize}

	\medskip
	\textbf{\textcolor{CExplanation}{特别提醒}}

	仅知道 $\sin x$ 或仅知道 $\cos x$ 时，$\cos 2x$ 往往可以由平方关系确定，但 $\sin 2x=2\sin x\cos x$ 可能因另一个三角函数的符号不确定而不能唯一确定。

\end{explanationbox}